\documentclass[11pt]{article}
\usepackage[margin=2.2cm]{geometry}
\usepackage[T1]{fontenc}
\usepackage{lmodern}
\usepackage{amsmath}
\usepackage{booktabs}
\usepackage{longtable}
\usepackage{multicol}
\usepackage{xcolor}
\usepackage{hyperref}

\newcommand{\todo}[1]{\par\noindent\fcolorbox{red}{yellow!25}{\parbox{\dimexpr\linewidth-2\fboxsep-2\fboxrule}{\textbf{TODO (by hand):} #1}}\par\smallskip}
\newcommand{\G}{\ensuremath{G_1}}

\title{Formal Languages Project\\FIRST/FOLLOW, LL(1), SLR(1) and SDD evaluation}
\author{Mateo Sepúlveda Escobar \and Isabella Hernández Posada}
\date{SI2002 Formal Languages --- EAFIT, 2026-2}

\begin{document}
\maketitle

\section{Overview}

We implemented, in Python using only its standard library, the algorithms of
Aho et al., \emph{Compilers: Principles, Techniques, and Tools} (2nd ed.)
required by the assignment. Table~\ref{tab:map} relates each part of the book
to the file that implements it and to the data structures it uses.

\begin{table}[h]
\centering\small
\begin{tabular}{p{4.6cm}p{2.2cm}p{8.2cm}}
\toprule
Book & File & Data structures \\
\midrule
FIRST and FOLLOW, Section 4.4.2 &
\texttt{grammar.py} &
Productions as a list of pairs \texttt{(head, body)}, where the body is a tuple of
symbols and $\varepsilon$ is the empty tuple. FIRST and FOLLOW are dictionaries
from symbol to set. \\
LL(1) grammars, Section 4.4.3 (definition and Algorithm 4.31) &
\texttt{ll1.py} &
The table is a dictionary $(A, a) \mapsto$ list of production indices. \\
Table-driven predictive parsing, Section 4.4.4 (Algorithm 4.34) &
\texttt{ll1.py} &
Stack of parse-tree nodes (a Python list; the top is the last element). \\
LR(0) items, CLOSURE, GOTO and canonical collection, Section 4.6.2 &
\texttt{slr.py} &
An item is a pair (production index, dot position). A state is a
\texttt{frozenset} of items. A dictionary from state to number avoids duplicates;
transitions are a dictionary $(i, X) \mapsto j$. \\
SLR table, Section 4.6.4 (Algorithm 4.46) &
\texttt{slr.py} &
ACTION: dictionary $(i, a) \mapsto$ shift $j$ / reduce $p$ / accept.
GOTO: dictionary $(i, A) \mapsto j$. \\
LR parsing, Section 4.6.3 (Algorithm 4.44) &
\texttt{slr.py} &
Stack of states and a parallel stack of parse-tree nodes. \\
SDDs, Sections 5.1.1--5.1.2 &
\texttt{sdd.py}, \texttt{sdd\_def.py} &
A rule stores the attribute it defines, the attributes it uses (as
(position, name) pairs, position 0 being the head) and a function. \\
Dependency graphs and evaluation order, Sections 5.2.1--5.2.2 &
\texttt{sdd.py} &
Vertices are pairs (node id, attribute); a dictionary maps each vertex to the set
of vertices it depends on. Topological order with Kahn's algorithm and a heap. \\
\bottomrule
\end{tabular}
\caption{Book sections, files and data structures.}
\label{tab:map}
\end{table}

The programs are \texttt{main.py} (Items 1 and 2, with the input and output of
Section 2.2 of the statement) and \texttt{sdd\_main.py} (Item 3). Both parsers
offer the same interface (\texttt{name}, \texttt{is\_valid}, \texttt{parse}),
so the program can choose one at run time and the SDD evaluator can use the
parse tree built by either of them. Instructions to run them are in
\texttt{README.md}.

\section{Items 1 and 2: FIRST, FOLLOW, LL(1) and SLR(1)}

\subsection{Algorithms}

\paragraph{FIRST and FOLLOW (Section 4.4.2).}
FIRST$(a) = \{a\}$ for every terminal. For the nonterminals we apply the rules
of the book until no set changes: for each production $A \to X_1 \cdots X_n$
we add FIRST$(X_1 \cdots X_n)$ to FIRST$(A)$. FIRST of a string adds
FIRST$(X_i) - \{\varepsilon\}$ while $X_1, \ldots, X_{i-1}$ are all nullable,
and contains $\varepsilon$ only if every $X_i$ is nullable (in particular, for
the empty string). FOLLOW starts with $\$ \in$ FOLLOW$(S)$; for each
$A \to \alpha B \beta$ we add FIRST$(\beta) - \{\varepsilon\}$ to FOLLOW$(B)$,
and FOLLOW$(A)$ when $\varepsilon \in$ FIRST$(\beta)$, again until nothing
changes.

\paragraph{LL(1) test (Section 4.4.3).}
We implement the definition at the beginning of the section and do not
transform the grammar. For every pair of distinct productions
$A \to \alpha \mid \beta$: (1) and (2) hold iff
FIRST$(\alpha) \cap$ FIRST$(\beta) = \emptyset$ (with $\varepsilon$ counted
as an element); (3) if $\varepsilon \in$ FIRST$(\beta)$ then
FIRST$(\alpha) \cap$ FOLLOW$(A) = \emptyset$, and symmetrically.

\paragraph{Predictive table and parser (Algorithms 4.31 and 4.34).}
For each $A \to \alpha$ we add the production to $M[A, a]$ for every
$a \in$ FIRST$(\alpha)$, and to $M[A, b]$ for every $b \in$ FOLLOW$(A)$
(including \$) if $\varepsilon \in$ FIRST$(\alpha)$. The parser starts with
$S\$$ on the stack and $w\$$ as input; it matches terminals, expands
nonterminals with $M[X, a]$, and reports an error on an empty entry. Each
expansion creates the children of the node on top of the stack, so the parse
tree is built while parsing.

\paragraph{SLR(1) (Sections 4.5, 4.6.2--4.6.4).}
We augment the grammar with production 0, $S' \to S$. CLOSURE adds
$B \to \cdot\gamma$ for each item with the dot before a nonterminal $B$;
GOTO$(I, X)$ moves the dot over $X$ and takes the closure. The canonical
collection starts with CLOSURE$(\{[S' \to \cdot S]\})$ and adds GOTO$(I, X)$
for every state and symbol until no new state appears. ACTION gets
shift $j$ for $[A \to \alpha \cdot a \beta]$, reduce $A \to \alpha$ for every
$a \in$ FOLLOW$(A)$ for $[A \to \alpha \cdot]$, and accept for
$[S' \to S \cdot]$ on \$. A second, different action in the same entry is a
conflict, and the grammar is SLR(1) iff there is none. The parser of
Algorithm 4.44 builds the tree bottom-up: a reduction by $A \to \beta$ pops
$|\beta|$ nodes and makes them the children of a new node $A$.

\subsection{Our grammar \G}

\G\ describes records: a record is a list in parentheses of fields
\texttt{a:V} separated by commas, and a value is \texttt{0}, \texttt{1} or
another record. It has $\varepsilon$-productions and is recursive.

\begin{verbatim}
S -> (P)
P -> KM e
M -> ,KM e
K -> a:V
V -> 0 1 S
\end{verbatim}

\noindent The numbering used by the SLR table (\texttt{r}$p$ means reduce by
production $p$) is:

\begin{center}\small
\begin{tabular}{rl}
0 & \texttt{S'} $\to$ \texttt{S} \\
1 & \texttt{S} $\to$ \texttt{(P)} \\
2 & \texttt{P} $\to$ \texttt{KM} \\
3 & \texttt{P} $\to$ $\varepsilon$ \\
4 & \texttt{M} $\to$ \texttt{,{}KM} \\
5 & \texttt{M} $\to$ $\varepsilon$ \\
6 & \texttt{K} $\to$ \texttt{a:V} \\
7 & \texttt{V} $\to$ \texttt{0} \\
8 & \texttt{V} $\to$ \texttt{1} \\
9 & \texttt{V} $\to$ \texttt{S}
\end{tabular}
\end{center}

\noindent FIRST and FOLLOW, as computed by our program (\texttt{main.py -{}-sets}):

\begin{center}\small
\begin{tabular}{lll}
\toprule
Nonterminal & FIRST & FOLLOW \\
\midrule
\texttt{K} & \{\texttt{a}\} & \{\texttt{)}, \texttt{,{}}\} \\
\texttt{M} & \{\texttt{,{}}, $\varepsilon$\} & \{\texttt{)}\} \\
\texttt{P} & \{\texttt{a}, $\varepsilon$\} & \{\texttt{)}\} \\
\texttt{S} & \{\texttt{(}\} & \{\texttt{)}, \texttt{,{}}, \texttt{\$}\} \\
\texttt{V} & \{\texttt{(}, \texttt{0}, \texttt{1}\} & \{\texttt{)}, \texttt{,{}}\} \\
\bottomrule
\end{tabular}
\end{center}

\noindent \G\ is LL(1) and SLR(1). The LL(1) table computed by our program
(\texttt{main.py -{}-tables}); empty entries are errors:

\begin{center}\small
\begin{tabular}{lcccccccc}
\toprule
 & \texttt{(} & \texttt{)} & \texttt{,{}} & \texttt{0} & \texttt{1} & \texttt{:} & \texttt{a} & \texttt{\$} \\
\midrule
\texttt{K} &  &  &  &  &  &  & \texttt{K} $\to$ \texttt{a:V} &  \\
\texttt{M} &  & \texttt{M} $\to$ $\varepsilon$ & \texttt{M} $\to$ \texttt{,{}KM} &  &  &  &  &  \\
\texttt{P} &  & \texttt{P} $\to$ $\varepsilon$ &  &  &  &  & \texttt{P} $\to$ \texttt{KM} &  \\
\texttt{S} & \texttt{S} $\to$ \texttt{(P)} &  &  &  &  &  &  &  \\
\texttt{V} & \texttt{V} $\to$ \texttt{S} &  &  & \texttt{V} $\to$ \texttt{0} & \texttt{V} $\to$ \texttt{1} &  &  &  \\
\bottomrule
\end{tabular}
\end{center}

\noindent The canonical collection of LR(0) items (kernel items first):

\begin{multicols}{4}\small
\noindent\textbf{I0}\\\texttt{S\textquotesingle} $\to$ \texttt{$\cdot$S}\\\texttt{S} $\to$ \texttt{$\cdot$(P)}

\smallskip
\noindent\textbf{I1}\\\texttt{S\textquotesingle} $\to$ \texttt{S$\cdot$}

\smallskip
\noindent\textbf{I2}\\\texttt{S} $\to$ \texttt{($\cdot$P)}\\\texttt{P} $\to$ \texttt{$\cdot$KM}\\\texttt{P} $\to$ \texttt{$\cdot$}\\\texttt{K} $\to$ \texttt{$\cdot$a:V}

\smallskip
\noindent\textbf{I3}\\\texttt{P} $\to$ \texttt{K$\cdot$M}\\\texttt{M} $\to$ \texttt{$\cdot$,{}KM}\\\texttt{M} $\to$ \texttt{$\cdot$}

\smallskip
\noindent\textbf{I4}\\\texttt{S} $\to$ \texttt{(P$\cdot$)}

\smallskip
\noindent\textbf{I5}\\\texttt{K} $\to$ \texttt{a$\cdot$:V}

\smallskip
\noindent\textbf{I6}\\\texttt{P} $\to$ \texttt{KM$\cdot$}

\smallskip
\noindent\textbf{I7}\\\texttt{M} $\to$ \texttt{,{}$\cdot$KM}\\\texttt{K} $\to$ \texttt{$\cdot$a:V}

\smallskip
\noindent\textbf{I8}\\\texttt{S} $\to$ \texttt{(P)$\cdot$}

\smallskip
\noindent\textbf{I9}\\\texttt{K} $\to$ \texttt{a:$\cdot$V}\\\texttt{S} $\to$ \texttt{$\cdot$(P)}\\\texttt{V} $\to$ \texttt{$\cdot$0}\\\texttt{V} $\to$ \texttt{$\cdot$1}\\\texttt{V} $\to$ \texttt{$\cdot$S}

\smallskip
\noindent\textbf{I10}\\\texttt{M} $\to$ \texttt{,{}K$\cdot$M}\\\texttt{M} $\to$ \texttt{$\cdot$,{}KM}\\\texttt{M} $\to$ \texttt{$\cdot$}

\smallskip
\noindent\textbf{I11}\\\texttt{V} $\to$ \texttt{S$\cdot$}

\smallskip
\noindent\textbf{I12}\\\texttt{K} $\to$ \texttt{a:V$\cdot$}

\smallskip
\noindent\textbf{I13}\\\texttt{V} $\to$ \texttt{0$\cdot$}

\smallskip
\noindent\textbf{I14}\\\texttt{V} $\to$ \texttt{1$\cdot$}

\smallskip
\noindent\textbf{I15}\\\texttt{M} $\to$ \texttt{,{}KM$\cdot$}
\end{multicols}

\noindent Transitions: {\small $I_{0}\xrightarrow{\texttt{(}}I_{2}$, $I_{0}\xrightarrow{\texttt{S}}I_{1}$, $I_{2}\xrightarrow{\texttt{K}}I_{3}$, $I_{2}\xrightarrow{\texttt{P}}I_{4}$, $I_{2}\xrightarrow{\texttt{a}}I_{5}$, $I_{3}\xrightarrow{\texttt{,{}}}I_{7}$, $I_{3}\xrightarrow{\texttt{M}}I_{6}$, $I_{4}\xrightarrow{\texttt{)}}I_{8}$, $I_{5}\xrightarrow{\texttt{:}}I_{9}$, $I_{7}\xrightarrow{\texttt{K}}I_{10}$, $I_{7}\xrightarrow{\texttt{a}}I_{5}$, $I_{9}\xrightarrow{\texttt{(}}I_{2}$, $I_{9}\xrightarrow{\texttt{0}}I_{13}$, $I_{9}\xrightarrow{\texttt{1}}I_{14}$, $I_{9}\xrightarrow{\texttt{S}}I_{11}$, $I_{9}\xrightarrow{\texttt{V}}I_{12}$, $I_{10}\xrightarrow{\texttt{,{}}}I_{7}$, $I_{10}\xrightarrow{\texttt{M}}I_{15}$.}

\begin{center}\small
\begin{tabular}{r|cccccccc|ccccc}
\toprule
state & \texttt{(} & \texttt{)} & \texttt{,{}} & \texttt{0} & \texttt{1} & \texttt{:} & \texttt{a} & \texttt{\$} & \texttt{K} & \texttt{M} & \texttt{P} & \texttt{S} & \texttt{V} \\
\midrule
0 & s2 &  &  &  &  &  &  &  &  &  &  & 1 &  \\
1 &  &  &  &  &  &  &  & acc &  &  &  &  &  \\
2 &  & r3 &  &  &  &  & s5 &  & 3 &  & 4 &  &  \\
3 &  & r5 & s7 &  &  &  &  &  &  & 6 &  &  &  \\
4 &  & s8 &  &  &  &  &  &  &  &  &  &  &  \\
5 &  &  &  &  &  & s9 &  &  &  &  &  &  &  \\
6 &  & r2 &  &  &  &  &  &  &  &  &  &  &  \\
7 &  &  &  &  &  &  & s5 &  & 10 &  &  &  &  \\
8 &  & r1 & r1 &  &  &  &  & r1 &  &  &  &  &  \\
9 & s2 &  &  & s13 & s14 &  &  &  &  &  &  & 11 & 12 \\
10 &  & r5 & s7 &  &  &  &  &  &  & 15 &  &  &  \\
11 &  & r9 & r9 &  &  &  &  &  &  &  &  &  &  \\
12 &  & r6 & r6 &  &  &  &  &  &  &  &  &  &  \\
13 &  & r7 & r7 &  &  &  &  &  &  &  &  &  &  \\
14 &  & r8 & r8 &  &  &  &  &  &  &  &  &  &  \\
15 &  & r4 &  &  &  &  &  &  &  &  &  &  &  \\
\bottomrule
\end{tabular}
\end{center}

\subsection{A string parsed step by step}

Trace of \texttt{(a:1,{}a:())} with the LL(1) parser (\texttt{main.py -{}-trace});
the top of the stack is on the left, as in Fig.~4.21 of the book:

{\small
\begin{longtable}{rlll}
\toprule
& Stack & Input & Action \\
\midrule
1 & \texttt{S\$} & \texttt{(a:1,{}a:())\$} & \texttt{output S $\to$ (P)} \\
2 & \texttt{(P)\$} & \texttt{(a:1,{}a:())\$} & \texttt{match (} \\
3 & \texttt{P)\$} & \texttt{a:1,{}a:())\$} & \texttt{output P $\to$ KM} \\
4 & \texttt{KM)\$} & \texttt{a:1,{}a:())\$} & \texttt{output K $\to$ a:V} \\
5 & \texttt{a:VM)\$} & \texttt{a:1,{}a:())\$} & \texttt{match a} \\
6 & \texttt{:VM)\$} & \texttt{:1,{}a:())\$} & \texttt{match :} \\
7 & \texttt{VM)\$} & \texttt{1,{}a:())\$} & \texttt{output V $\to$ 1} \\
8 & \texttt{1M)\$} & \texttt{1,{}a:())\$} & \texttt{match 1} \\
9 & \texttt{M)\$} & \texttt{,{}a:())\$} & \texttt{output M $\to$ ,{}KM} \\
10 & \texttt{,{}KM)\$} & \texttt{,{}a:())\$} & \texttt{match ,{}} \\
11 & \texttt{KM)\$} & \texttt{a:())\$} & \texttt{output K $\to$ a:V} \\
12 & \texttt{a:VM)\$} & \texttt{a:())\$} & \texttt{match a} \\
13 & \texttt{:VM)\$} & \texttt{:())\$} & \texttt{match :} \\
14 & \texttt{VM)\$} & \texttt{())\$} & \texttt{output V $\to$ S} \\
15 & \texttt{SM)\$} & \texttt{())\$} & \texttt{output S $\to$ (P)} \\
16 & \texttt{(P)M)\$} & \texttt{())\$} & \texttt{match (} \\
17 & \texttt{P)M)\$} & \texttt{))\$} & \texttt{output P $\to$ e} \\
18 & \texttt{)M)\$} & \texttt{))\$} & \texttt{match )} \\
19 & \texttt{M)\$} & \texttt{)\$} & \texttt{output M $\to$ e} \\
20 & \texttt{)\$} & \texttt{)\$} & \texttt{match )} \\
21 & \texttt{\$} & \texttt{\$} & \texttt{accept} \\
\bottomrule
\end{longtable}}

\noindent The same string with the SLR(1) parser (states and symbols on the
stack, top on the right):

{\small
\begin{longtable}{rlll}
\toprule
& Stack & Input & Action \\
\midrule
1 & \texttt{0} & \texttt{(a:1,{}a:())\$} & \texttt{shift 2} \\
2 & \texttt{0 ( 2} & \texttt{a:1,{}a:())\$} & \texttt{shift 5} \\
3 & \texttt{0 ( 2 a 5} & \texttt{:1,{}a:())\$} & \texttt{shift 9} \\
4 & \texttt{0 ( 2 a 5 : 9} & \texttt{1,{}a:())\$} & \texttt{shift 14} \\
5 & \texttt{0 ( 2 a 5 : 9 1 14} & \texttt{,{}a:())\$} & \texttt{reduce V $\to$ 1} \\
6 & \texttt{0 ( 2 a 5 : 9 V 12} & \texttt{,{}a:())\$} & \texttt{reduce K $\to$ a:V} \\
7 & \texttt{0 ( 2 K 3} & \texttt{,{}a:())\$} & \texttt{shift 7} \\
8 & \texttt{0 ( 2 K 3 ,{} 7} & \texttt{a:())\$} & \texttt{shift 5} \\
9 & \texttt{0 ( 2 K 3 ,{} 7 a 5} & \texttt{:())\$} & \texttt{shift 9} \\
10 & \texttt{0 ( 2 K 3 ,{} 7 a 5 : 9} & \texttt{())\$} & \texttt{shift 2} \\
11 & \texttt{0 ( 2 K 3 ,{} 7 a 5 : 9 ( 2} & \texttt{))\$} & \texttt{reduce P $\to$ e} \\
12 & \texttt{0 ( 2 K 3 ,{} 7 a 5 : 9 ( 2 P 4} & \texttt{))\$} & \texttt{shift 8} \\
13 & \texttt{0 ( 2 K 3 ,{} 7 a 5 : 9 ( 2 P 4 ) 8} & \texttt{)\$} & \texttt{reduce S $\to$ (P)} \\
14 & \texttt{0 ( 2 K 3 ,{} 7 a 5 : 9 S 11} & \texttt{)\$} & \texttt{reduce V $\to$ S} \\
15 & \texttt{0 ( 2 K 3 ,{} 7 a 5 : 9 V 12} & \texttt{)\$} & \texttt{reduce K $\to$ a:V} \\
16 & \texttt{0 ( 2 K 3 ,{} 7 K 10} & \texttt{)\$} & \texttt{reduce M $\to$ e} \\
17 & \texttt{0 ( 2 K 3 ,{} 7 K 10 M 15} & \texttt{)\$} & \texttt{reduce M $\to$ ,{}KM} \\
18 & \texttt{0 ( 2 K 3 M 6} & \texttt{)\$} & \texttt{reduce P $\to$ KM} \\
19 & \texttt{0 ( 2 P 4} & \texttt{)\$} & \texttt{shift 8} \\
20 & \texttt{0 ( 2 P 4 ) 8} & \texttt{\$} & \texttt{reduce S $\to$ (P)} \\
21 & \texttt{0 S 1} & \texttt{\$} & \texttt{accept} \\
\bottomrule
\end{longtable}}

\subsection{Results}

\begin{center}\small
\begin{tabular}{lcc}
\toprule
String & LL(1) & SLR(1) \\
\midrule
\texttt{()} & yes & yes \\
\texttt{(a:1,{}a:(a:0))} & yes & yes \\
\texttt{(a:0,{}a:1,{}a:(a:(a:1)))} & yes & yes \\
\texttt{(a:1,{})} & no & no \\
\texttt{(a:(),{}a:1} & no & no \\
\texttt{e} & no & no \\
\texttt{(b:1)} & no & no \\
\bottomrule
\end{tabular}
\end{center}

\noindent \texttt{e} is the empty string, which is not in $L(\G)$;
\texttt{b} is not a terminal of \G. The exact run of \texttt{main.py}:

{\small
\noindent Input:
\begin{multicols}{3}
\begin{verbatim}
5
S -> (P)
P -> KM e
M -> ,KM e
K -> a:V
V -> 0 1 S
T
()
(a:1,a:(a:0))
(a:0,a:1,a:(a:(a:1)))
(a:1,)
(a:(),a:1
e
(b:1)

B
()
(a:1,a:(a:0))
(a:0,a:1,a:(a:(a:1)))
(a:1,)
(a:(),a:1
e
(b:1)

Q
\end{verbatim}
\end{multicols}
\noindent Output:
\begin{verbatim}
Select a parser (T: for LL(1), B: for SLR(1), Q: quit):
yes
yes
yes
no
no
no
no
Select a parser (T: for LL(1), B: for SLR(1), Q: quit):
yes
yes
yes
no
no
no
no
Select a parser (T: for LL(1), B: for SLR(1), Q: quit):
\end{verbatim}}

\section{Item 3: syntax-directed definition}

\subsection{Grammar}

The language is the nonempty lists of digits separated by commas, where an
element can also be a nested list in brackets, for example
\texttt{1,[2,3],4}.

\begin{center}
\begin{tabular}{l}
\texttt{S -> L} \\
\texttt{L -> EM} \\
\texttt{M -> ,EM e} \\
\texttt{E -> D [L]} \\
\texttt{D -> 0 1 2 3 4 5 6 7 8 9}
\end{tabular}
\end{center}

It satisfies the conditions of Step 1. G1: it follows the conventions of
Section 2.2 and it is LL(1) and SLR(1); we parse with our own parsers (LL(1)
by default, SLR(1) with \texttt{-{}-slr}, and both give the same results).
G2: it has 5 nonterminals and 16 productions, and $L$ and $M$ are recursive
($L \Rightarrow EM \Rightarrow [L]M$ and $M \to ,EM$). G3: neither the grammar
nor the SDD is in the book or the statement. A first version of our SDD also
computed the maximum nesting depth; we removed that attribute because it was
the same idea as the example of the statement.

\subsection{Attributes}

The SDD computes the \emph{depth-weighted sum} of the list: every digit is
multiplied by the depth of the list where it appears (the top level is depth
1). For \texttt{1,[2,3],4} it is $1\cdot1 + (2+3)\cdot2 + 4\cdot1 = 15$.

\begin{center}\small
\begin{tabular}{lllll}
\toprule
Symbol & Attribute & Kind & Type & Meaning \\
\midrule
$S$ & $s$ & synthesized & integer & depth-weighted sum of the whole list (main attribute) \\
$L$, $M$, $E$ & $d$ & inherited & integer & depth of the list the symbol belongs to \\
$L$, $M$, $E$ & $s$ & synthesized & integer & depth-weighted sum of the digits in the subtree \\
$D$ & $val$ & synthesized & integer & numeric value of the digit \\
terminals & $lexval$ & synthesized & character & the character itself (given by the scanner) \\
\bottomrule
\end{tabular}
\end{center}

\subsection{Semantic rules}

In the format of Fig.~5.1 of the book (subscripts distinguish two occurrences
of the same nonterminal):

\begin{center}
\begin{tabular}{ll}
\toprule
Production & Semantic rules \\
\midrule
$S \to L$ & $L.d = 1$; \quad $S.s = L.s$ \\
$L \to E\,M$ & $E.d = L.d$; \quad $M.d = L.d$; \quad $L.s = E.s + M.s$ \\
$M \to ,\,E\,M_1$ & $E.d = M.d$; \quad $M_1.d = M.d$; \quad $M.s = E.s + M_1.s$ \\
$M \to \varepsilon$ & $M.s = 0$ \\
$E \to D$ & $E.s = D.val \times E.d$ \\
$E \to [\,L\,]$ & $L.d = E.d + 1$; \quad $E.s = L.s$ \\
$D \to \mathit{digit}$ & $D.val = \mathit{digit}.lexval$ as an integer \quad ($\mathit{digit} \in \{0, \ldots, 9\}$) \\
\bottomrule
\end{tabular}
\end{center}

We checked the conditions of Section 5.1.1 for every production, and the
program checks them too before evaluating (\texttt{SDD.check}): each rule
defines a synthesized attribute of the head or an inherited attribute of a
symbol of the body, using only attributes of the same production; $s$ of the
head and $d$ of every nonterminal of the body are defined exactly once; and
$S$ has no inherited attributes.

\subsection{Example by hand (Step 3)}

We use the string \texttt{1,[2,3]} (7 terminals).

\todo{Draw the parse tree of \texttt{1,[2,3]} and annotate every node with the
values of its attributes.}

\todo{Draw the dependency graph of that tree: one vertex for each attribute of
each node, and one edge from each attribute that a rule uses to the attribute
that the rule defines (Section 5.2.1).}

\todo{Give one valid evaluation order (a topological order of the graph,
Section 5.2.2).}

\noindent We verified the example with our program. Nodes are numbered in
preorder, and \texttt{X\#n.a} is attribute \texttt{a} of node \texttt{n}:

{\small
\begin{multicols}{2}
\begin{verbatim}
Annotated parse tree:
S#1 [S -> L] {s=11}
  L#2 [L -> EM] {d=1, s=11}
    E#3 [E -> D] {d=1, s=1}
      D#4 [D -> 1] {val=1}
        '1'#5 {lexval=1}
    M#6 [M -> ,EM] {d=1, s=10}
      ','#7 {lexval=,}
      E#8 [E -> [L]] {d=1, s=10}
        '['#9 {lexval=[}
        L#10 [L -> EM] {d=2, s=10}
          E#11 [E -> D] {d=2, s=4}
            D#12 [D -> 2] {val=2}
              '2'#13 {lexval=2}
          M#14 [M -> ,EM] {d=2, s=6}
            ','#15 {lexval=,}
            E#16 [E -> D] {d=2, s=6}
              D#17 [D -> 3] {val=3}
                '3'#18 {lexval=3}
            M#19 [M -> e] {d=2, s=0}
        ']'#20 {lexval=]}
      M#21 [M -> e] {d=1, s=0}
Evaluation order:
    1. L#2.d
    2. E#3.d
    3. '1'#5.lexval   (given by the scanner)
    4. D#4.val
    5. E#3.s
    6. M#6.d
    7. ','#7.lexval   (given by the scanner)
    8. E#8.d
    9. '['#9.lexval   (given by the scanner)
   10. L#10.d
   11. E#11.d
   12. '2'#13.lexval   (given by the scanner)
   13. D#12.val
   14. E#11.s
   15. M#14.d
   16. ','#15.lexval   (given by the scanner)
   17. E#16.d
   18. '3'#18.lexval   (given by the scanner)
   19. D#17.val
   20. E#16.s
   21. M#19.d
   22. M#19.s
   23. M#14.s
   24. L#10.s
   25. E#8.s
   26. ']'#20.lexval   (given by the scanner)
   27. M#21.d
   28. M#21.s
   29. M#6.s
   30. L#2.s
   31. S#1.s
S.s = 11
\end{verbatim}
\end{multicols}}

\subsection{Classification (Step 4)}

The SDD is \textbf{not S-attributed}, because $d$ is an inherited attribute
(Section 5.2.3 requires every attribute to be synthesized).

It is \textbf{L-attributed} (Section 5.2.4). Synthesized attributes are always
allowed. Every inherited attribute of a body symbol $X_i$ depends only on
inherited attributes of the head and on attributes of symbols to the left of
$X_i$: $L.d = 1$ uses nothing; $E.d = L.d$, $M.d = L.d$, $E.d = M.d$,
$M_1.d = M.d$ and $L.d = E.d + 1$ use only the inherited attribute $d$ of the
head. No inherited attribute depends on an attribute of a symbol to its right,
and no attribute of a node depends on itself.

As the statement requires, our evaluator does not use this property: it
always builds the dependency graph and follows a topological order, so it
works for any non-circular SDD.

\subsection{Implementation (Step 5)}

\begin{enumerate}
\item \emph{SDD as data.} \texttt{sdd\_def.py} contains the grammar lines, the
attributes of each symbol and, for each production, its list of rules. A rule
is (attribute defined, attributes used, function), with attributes written as
(position, name). \texttt{sdd.py} does not know anything about our SDD.
\item \emph{Parse tree.} The parsers of Item 2 build it. Each node stores its
symbol, the index of the production used (nonterminals only), its children
and its attribute values; leaves get \textit{lexval}.
\item \emph{Dependency graph.} For each internal node we instantiate the rules
of its production, replacing each position by the corresponding node, and
add one edge from each attribute used to the attribute defined.
\item \emph{Order and evaluation.} Kahn's algorithm gives a topological order
(ties are broken by node number to make the output deterministic). If some
vertices are left, the graph has a cycle: we report it and stop. Otherwise we
apply the rules in that order.
\item \emph{Output.} For each string, \texttt{no} if it is not in the
language; otherwise the annotated tree (one node per line, indented by depth),
the evaluation order and the value of $S.s$.
\end{enumerate}

\subsection{Tests (Step 6)}

Results with \texttt{sdd\_main.py}; the LL(1) and the SLR(1) parsers give
the same values:

\begin{center}\small
\begin{tabular}{lll}
\toprule
String & Output & Note \\
\midrule
\texttt{7} & \texttt{S.s = 7} & shortest string (one digit) \\
\texttt{1,{}[2,{}3],{}4} & \texttt{S.s = 15} &  \\
\texttt{[1,{}[2]],{}3} & \texttt{S.s = 11} &  \\
\texttt{[[[[5]]]]} & \texttt{S.s = 25} &  \\
\texttt{[}$^{30}$\texttt{1}\texttt{]}$^{30}$ & \texttt{S.s = 31} & deep nesting: 30 levels \\
\texttt{1,{},{}2} & \texttt{no} & not in the language \\
\texttt{[1} & \texttt{no} & not in the language \\
\bottomrule
\end{tabular}
\end{center}

\noindent To test circularity, \texttt{sdd\_main.py -{}-circular} replaces the
rule $L.d = 1$ of $S \to L$ by $L.d = L.s$. Since $L.s$ depends on $E.s$,
which depends on $E.d$, which depends on $L.d$, the graph has a cycle and the
program stops:

{\small
\begin{verbatim}
The SDD is circular. Cycle in the dependency graph:
  L#2.d -> E#3.d -> E#3.s -> L#2.s -> L#2.d
\end{verbatim}}

\section{Testing}

Besides the examples above, \texttt{tests/test\_project.py} checks
automatically: the three examples of the statement (exact output, including
\texttt{-{}-sets}); edge cases (terminals such as \texttt{-}, \texttt{\textgreater} and
\texttt{|}, the empty string, the line of $S$ not being the first one,
symbols that are not terminals, end of input without an empty line); a string
of 100\,000 symbols; and the SDD with both parsers, a nesting depth of 2000
and the circular variant. It also generates hundreds of random grammars and
compares the answers of both parsers, for all strings of length up to 4, with
an independent Earley recognizer.

\section{Problems found and how we solved them}

\begin{enumerate}
\item \textbf{LL(1) test vs.\ table conflicts.} We first decided LL(1)-ness by
looking for multiply-defined entries in the table. The statement asks for the
definition of Section 4.4.3, so we implemented it. Comparing both on random
grammars showed they differ in one case: an unreachable nonterminal with two
nullable alternatives (e.g.\ \texttt{S -> a e}, \texttt{B -> S e}). FOLLOW$(B)$
is empty, so the table has no conflict, but condition (2) fails. We follow the
definition, so our program says that grammar is SLR(1) only.
\item \textbf{Reductions by $\varepsilon$.} A reduction by $A \to \varepsilon$
pops zero symbols. In Python, \texttt{nodes[-0:]} is the whole list, not an
empty one, so we slice with \texttt{nodes[len(nodes) - k:]}.
\item \textbf{Augmented start symbol.} The new start symbol must not be a
symbol of the grammar. Since all symbols are single characters, we use the
two-character name \texttt{S'}.
\item \textbf{Input strings.} The end marker is appended by the program, and
a string with a symbol that is not a terminal of the grammar (for example
\texttt{\$}) is rejected before parsing; otherwise \texttt{a\$} could be
accepted by mistake.
\item \textbf{Slow parsing of long strings.} Our first LL(1) parser built the
text of the trace in every step even when it was not printed, which made it
quadratic: a string of 32\,000 symbols took more than 3 minutes. Now the trace
is built only with \texttt{-{}-trace}, and the same string takes less than one
second.
\item \textbf{Deep trees.} Recursive traversals of the parse tree failed with
deep nesting (Python's recursion limit). We made them iterative.
\item \textbf{Interactive use.} The first version read all the input before
answering, so the menu of case 1 did not respond when typing. Now the program
reads one line at a time.
\item \textbf{Module names.} Files named \texttt{sdd-main.py} or \texttt{LL1.py}
could not be imported as \texttt{sdd\_main} or \texttt{ll1}; we renamed them.
\end{enumerate}

\end{document}
